\chapter{Практика 8. Расчет BER и построение графиков}

\section{Задание}
Реализовать расчет вероятности битовых ошибок (BER) и построение графиков спектров и сигнальных созвездий.

\section{Теоретическое описание}
Вероятность битовой ошибки (Bit Error Rate, BER) вычисляется по формуле:
\[
BER = \frac{\text{количество ошибочных бит}}{\text{общее количество бит}}
\]

Сравнение производится между битовым сообщением после знакового кодера (на передатчике) и перед знаковым декодером (в приемнике).

\section{Реализация}

\subsection{Расчет BER}

\begin{lstlisting}[language=Python, caption=Расчет BER]
def calculate_ber(tx_bits, rx_bits):
    if len(tx_bits) != len(rx_bits):
        min_len = min(len(tx_bits), len(rx_bits))
        tx_bits = tx_bits[:min_len]
        rx_bits = rx_bits[:min_len]
    errors = sum(1 for i in range(len(tx_bits)) if tx_bits[i] != rx_bits[i])
    ber = errors / len(tx_bits) if len(tx_bits) > 0 else 0
    return ber, errors
\end{lstlisting}

\subsection{Построение графиков}

\begin{lstlisting}[language=Python, caption=Построение графиков спектров]
def plot_spectrums(tx_spectrum, rx_spectrum_before, rx_spectrum_after):
    plt.figure(figsize=(12, 8))
    
    plt.subplot(3, 1, 1)
    plt.plot(np.abs(tx_spectrum))
    plt.title('Спектр переданного OFDM символа')
    plt.xlabel('Индекс поднесущей')
    plt.ylabel('Амплитуда')
    plt.grid(True)
    
    plt.subplot(3, 1, 2)
    plt.plot(np.abs(rx_spectrum_before))
    plt.title('Спектр принятого OFDM символа до эквалайзирования')
    plt.xlabel('Индекс поднесущей')
    plt.ylabel('Амплитуда')
    plt.grid(True)
    
    plt.subplot(3, 1, 3)
    plt.plot(np.abs(rx_spectrum_after))
    plt.title('Спектр принятого OFDM символа после эквалайзирования')
    plt.xlabel('Индекс поднесущей')
    plt.ylabel('Амплитуда')
    plt.grid(True)
    
    plt.tight_layout()
    plt.show()
\end{lstlisting}

\begin{lstlisting}[language=Python, caption=Построение сигнальных созвездий]
def plot_constellations(tx_symbols, rx_symbols):
    plt.figure(figsize=(12, 5))
    
    plt.subplot(1, 2, 1)
    plt.scatter([s.real for s in tx_symbols], [s.imag for s in tx_symbols])
    plt.title('Сигнальное созвездие QPSK в передатчике')
    plt.xlabel('I')
    plt.ylabel('Q')
    plt.grid(True)
    plt.axis('equal')
    
    plt.subplot(1, 2, 2)
    plt.scatter([s.real for s in rx_symbols], [s.imag for s in rx_symbols])
    plt.title('Сигнальное созвездие QPSK в приемнике')
    plt.xlabel('I')
    plt.ylabel('Q')
    plt.grid(True)
    plt.axis('equal')
    
    plt.tight_layout()
    plt.show()
\end{lstlisting}

\section{Результат}
\begin{verbatim}
BER = 0.000000
Ошибочных бит: 0 из 408
\end{verbatim}

На графиках отображаются:
\begin{itemize}
    \item Спектр переданного OFDM символа
    \item Спектр принятого сигнала до эквалайзирования (с искажениями)
    \item Спектр принятого сигнала после эквалайзирования (восстановленный)
    \item Сигнальное созвездие QPSK в передатчике (4 точки)
    \item Сигнальное созвездие QPSK в приемнике (облака точек из-за шума)
\end{itemize}

\section{Вывод}
Расчет BER и построение графиков успешно реализованы. При заданных параметрах канала ($N_B = 3$, $N_0 = -100$ дБ) ошибок при передаче не обнаружено. Графики наглядно демонстрируют эффект эквалайзирования.